\documentclass[noauthor,noinstructornotes]{ximera}

\title{Week 1}
\author{Math 1193 Design Team}

\begin{document}
\begin{abstract}
    Activities for Week 1.
\end{abstract}
\maketitle

\begin{problem}
\begin{itemize}
    \item Factor the polynomial \(x^2 - x - 6\). Compare with 
your group and make sense of why each of your methods
works. If you all used the same method, try to find
a different one.
    \begin{itemize}
        \item Solve the equation \(x^2 - x - 6 = 0\).
    \end{itemize}
\end{itemize}

\begin{instructorNotes}
This is meant to draw out students' prior knowledge of factoring and expose them to others'.
\end{instructorNotes}
\end{problem}

\begin{problem}
\begin{itemize}
    \item Factor the polynomial \(x^2 - 25\). Compare with 
your group and make sense of why each of your methods
works. If you all used the same method, try to find
a different one.
    \begin{itemize}
        \item Solve the equation \(x^2 - 25 = 0\).
    \end{itemize}
\end{itemize}

\begin{instructorNotes}
This is meant to draw out students' prior knowledge of factoring and expose them to others'.
\end{instructorNotes}
\end{problem}

\begin{problem}
Solve the inequality \(-2x + 4 \ge 0\). Shade the 
solution set on the number line, then write the 
solution set in interval notation.

\begin{instructorNotes}
This is meant to see how students go about solving inequalities. 
Do they perform the operations in the correct order? Do they know 
about flipping the inequality? Also, it makes sure students understand 
intervals and interval notation.
\end{instructorNotes}
\end{problem}

\begin{problem}
Let's play a game (sort of)!

\begin{itemize}
    \item Consider the equation \(y = 2 + 3x\). Pair up and 
    choose one person to be Player X and the other
    to be Player Y.
    \begin{enumerate}
        \item Player X comes up with 3 \(x\)-values, and 
        Player Y comes up with 3 \(y\)-values. 
        These can be any numbers you want, but try 
        to keep them small so they can be plotted 
        later.
        \item The players exchange their numbers. 
        Player X has to find the \(x\)-values that 
        correspond to the \(y\)-values and put 
        them into coordinate pairs, and Player Y 
        has to find the \(y\)-values that 
        correspond to the \(x\)-values and put 
        them into coordinate pairs. 
        \begin{itemize}
            \item Whose job is easier?
        \end{itemize}
        \item Write down the 6 
        ordered pairs that satisfy the equation.
        \item Graph your ordered pairs on the \(xy\)
        coordinate plane.
    \end{enumerate}
\end{itemize}

\begin{instructorNotes}
Make sure to have paper at the students' tables and project
the problem or write on the board.
\end{instructorNotes}
\end{problem}

\begin{problem}
Jaden was given the following problem:

\textit{Let \(f(x) = x^2 + x + 1\). Find \(f(a + 2)\).}

Jaden's work is as follows:

\begin{align*}
    f(a + 2) & = a + 2^2 + a + 2 + 1 \\
    & = a + 4 + a + 3 \\
    & = 2a + 7
\end{align*}

Is Jaden's work correct? Why or why not? Where did he 
go wrong? What did he do correctly?

\begin{instructorNotes}
Print this out so students can annotate the work.
\end{instructorNotes}
\end{problem}

\begin{problem}
What are the domain and range of the function 
whose graph is shown below? Compare your answer with
a partner and discuss. 

\begin{center}
    \includegraphics[width=10cm]{domainRangeGraphically.png}
\end{center}

\begin{instructorNotes}
Can be printed so students can interact with the graph.
\end{instructorNotes}
\end{problem}

\end{document}
